%==============================================================================
%  CHAPTER 4 -- EXPERIMENTAL CAMPAIGN
%
%  Job of this chapter: let someone else repeat exactly what you did.
%  Past tense, impersonal. NO results here -- those go in Chapter 5.
%
%  The test: could a competent stranger rebuild your setup and rerun your
%  tests using only what is written here?
%==============================================================================
\chapter{Experimental campaign}
\label{ch:experimental}

%------------------------------------------------------------------------------
\section{Experimental test bench}
\label{sec:test-bench}
%  What the rig is, with a figure -- ideally a photograph AND a schematic.
%  List every instrument in a table: manufacturer, model, measurement range,
%  resolution and calibration status. In a measurements group, that table is
%  the first thing an examiner looks at.
%
%  \begin{table}[H]
%    \centering
%    \caption{Instrumentation used in the campaign.}
%    \label{tab:instruments}
%    \begin{tabular}{llll}
%      \toprule
%      Instrument & Manufacturer / model & Range & Resolution \\
%      \midrule
%      ... & ... & ... & ... \\
%      \bottomrule
%    \end{tabular}
%  \end{table}
%
%  If more than one bench was used, add a subsection for each.

%------------------------------------------------------------------------------
\section{Experimental protocol}
\label{sec:protocol}
%  The procedure, in the order performed: specimen preparation, mounting,
%  environmental conditions, excitation, acquisition settings (sampling rate,
%  duration, triggering) and the test matrix.
%
%  Say how many repetitions you ran, and why. "Three repetitions per
%  condition" is a result-bearing choice, not a detail.
%
%  State what you did NOT control, and what you assume it did not affect.

%------------------------------------------------------------------------------
\section{Data processing}
\label{sec:data-processing}
%  Everything done to the raw data before it becomes a result: filtering,
%  windowing, synchronisation, outlier rejection, averaging, unit conversion.
%
%  Give parameters, not just the names of steps. "A low-pass filter was
%  applied" is not reproducible. "A 4th-order zero-phase Butterworth low-pass
%  filter with a \SI{500}{\hertz} cut-off" is.
%
%  If you rejected data, say which, how much, and on what criterion -- decided
%  BEFORE looking at the results.
